\documentclass[../DoAn.tex]{subfiles}
\begin{document}
\section{Đặt vấn đề}
\label{section:1.1}

Hiện nay, thương mại điện tử và các kênh bán hàng trực tuyến đã trở thành một phần quan trọng trong hoạt động kinh doanh của các cửa hàng bán lẻ. Với lĩnh vực nội thất, việc bán hàng trực tuyến không chỉ là đăng hình ảnh, mô tả và giá sản phẩm. Khách hàng thường cần được tư vấn thêm về kích thước, chất liệu, phong cách, không gian sử dụng, ngân sách và mức độ phù hợp giữa nhiều lựa chọn khác nhau. Vì vậy, tư vấn bán hàng là một khâu quan trọng ảnh hưởng trực tiếp đến trải nghiệm và quyết định mua hàng của khách hàng.

Tuy nhiên, các cửa hàng nội thất nhỏ và vừa thường gặp khó khăn trong việc duy trì một quy trình tư vấn trực tuyến nhanh, nhất quán và đủ bám sát dữ liệu sản phẩm. Danh mục sản phẩm có thể nhiều và thay đổi theo thời gian, trong khi thông tin sản phẩm lại gồm nhiều thuộc tính như giá, chất liệu, kích thước, loại sản phẩm và phong cách. Nếu dữ liệu này không được tổ chức và khai thác tốt, việc tư vấn dễ phụ thuộc vào nhân viên, thiếu nhất quán giữa các lần trả lời và khó mở rộng khi số lượng khách hàng tăng.

Về phía người mua, nhu cầu khi chọn sản phẩm nội thất thường không cố định ngay từ đầu. Trong cùng một phiên tư vấn, khách hàng có thể đổi ngân sách, thay đổi mục đích sử dụng, yêu cầu gợi ý lại hoặc muốn so sánh nhiều sản phẩm trước khi quyết định. Đây là điểm yếu của các cách tư vấn đơn giản hoặc các chatbot chỉ xử lý theo kịch bản cố định. Từ điểm yếu đó có thể phát sinh nguy cơ chatbot đưa ra gợi ý không còn phù hợp với ngữ cảnh mới, khiến khách hàng mất thời gian, giảm niềm tin vào tư vấn trực tuyến và làm giảm khả năng chuyển đổi mua hàng.

Bên cạnh đó, khách hàng không chỉ cần một sản phẩm phù hợp mà còn muốn biết sản phẩm đó có đáng giá hay không so với các lựa chọn khác. Việc so sánh sản phẩm nội thất thường phải xét nhiều tiêu chí như giá, chất liệu, kích thước, kiểu dáng, phong cách và cửa hàng cung cấp. Nếu hệ thống chỉ hiển thị danh sách sản phẩm rời rạc, khách hàng phải tự đối chiếu thủ công, còn cửa hàng chưa khai thác được dữ liệu sẵn có để hỗ trợ quá trình tư vấn và tham chiếu giá.

Từ những vấn đề trên, đề tài hướng đến xây dựng hệ thống \textit{Multi-tenant RAG chatbot hỗ trợ tư vấn bán hàng, so sánh và tham chiếu giá sản phẩm nội thất}. Trong đó, \textit{multi-tenant} cho phép một hệ thống phục vụ nhiều cửa hàng, nhưng mỗi cửa hàng vẫn có dữ liệu sản phẩm, chatbot và cấu hình riêng. \textit{RAG} là hướng tiếp cận giúp chatbot truy xuất thông tin từ cơ sở tri thức trước khi tạo câu trả lời, nhờ đó nội dung tư vấn có thể bám sát dữ liệu sản phẩm hơn. Hệ thống được kỳ vọng hỗ trợ cửa hàng khai thác dữ liệu sản phẩm để tư vấn khách hàng hiệu quả hơn, đồng thời giúp người mua dễ dàng nhận gợi ý, so sánh sản phẩm và tham khảo mặt bằng giá trước khi ra quyết định.

\section{Mục tiêu và phạm vi đề tài}
\label{section:1.2}

Đề tài nhằm mục tiêu xây dựng một hệ thống chatbot hỗ trợ tư vấn bán hàng trong lĩnh vực nội thất. Hệ thống cho phép nhiều cửa hàng cùng sử dụng trên một nền tảng, nhưng mỗi cửa hàng vẫn có dữ liệu sản phẩm, chatbot và cấu hình riêng. Chatbot cần có khả năng tư vấn dựa trên dữ liệu sản phẩm của từng cửa hàng, duy trì ngữ cảnh hội thoại nhiều lượt và điều chỉnh gợi ý khi khách hàng thay đổi nhu cầu.

Cụ thể, đề tài tập trung vào các mục tiêu chính sau:

\begin{itemize}
    \item Xây dựng hệ thống theo mô hình multi-tenant, phục vụ nhiều cửa hàng nội thất trên cùng một nền tảng.
    \item Xây dựng chức năng quản lý dữ liệu sản phẩm hoặc tài liệu sản phẩm để tạo và cập nhật knowledge base cho từng cửa hàng.
    \item Xây dựng chatbot RAG có khả năng truy xuất thông tin từ knowledge base và tư vấn sản phẩm dựa trên dữ liệu liên quan.
    \item Hỗ trợ hội thoại nhiều lượt, trong đó chatbot có thể ghi nhận các thay đổi về ngân sách, nhu cầu, mục đích sử dụng hoặc yêu cầu gợi ý lại.
    \item Hỗ trợ so sánh nhiều sản phẩm theo các tiêu chí như giá, chất liệu, kích thước, loại sản phẩm, phong cách và cửa hàng cung cấp.
    \item Hỗ trợ tham chiếu giá dựa trên dữ liệu sản phẩm hiện có, giúp người dùng có thêm cơ sở khi đánh giá một sản phẩm hoặc một nhóm sản phẩm.
    \item Hỗ trợ tạo yêu cầu mua hàng khi người dùng đã có đủ thông tin cần thiết trong quá trình tư vấn.
\end{itemize}

Phạm vi của đề tài là xây dựng một hệ thống phần mềm demo gồm các thành phần chính: backend nghiệp vụ, AI service xử lý hội thoại và truy xuất tri thức, giao diện quản trị cho cửa hàng, giao diện quản trị hệ thống và giao diện chat cho người dùng cuối. Dữ liệu thử nghiệm gồm tối thiểu hai cửa hàng nội thất mẫu, trong đó mỗi cửa hàng có dữ liệu sản phẩm hoặc tài liệu sản phẩm riêng để phục vụ tư vấn.

Hệ thống tập trung vào ba nhóm chức năng chính. Nhóm thứ nhất là tư vấn theo từng cửa hàng, trong đó chatbot sử dụng dữ liệu riêng của cửa hàng để trả lời câu hỏi, gợi ý sản phẩm và tạo yêu cầu mua hàng. Nhóm thứ hai là tư vấn và so sánh chung trên dữ liệu tổng hợp, giúp người dùng tham khảo nhiều sản phẩm từ các cửa hàng khác nhau. Nhóm thứ ba là tham chiếu giá, trong đó hệ thống đưa ra thông tin về khoảng giá hoặc mặt bằng giá dựa trên dữ liệu hiện có.

Trong phạm vi đồ án, hệ thống không đặt mục tiêu xây dựng một sàn thương mại điện tử hoàn chỉnh. Các chức năng như thanh toán trực tuyến, quản lý vận chuyển, quản lý kho chuyên sâu hoặc chăm sóc sau bán hàng không phải là trọng tâm của đề tài. Chức năng tham chiếu giá cũng không nhằm phản ánh toàn bộ thị trường nội thất, mà chỉ cung cấp thông tin tham khảo dựa trên dữ liệu được đưa vào hệ thống.

\section{Định hướng giải pháp}
\label{section:1.3}

Từ các vấn đề đã nêu, đề tài định hướng xây dựng một hệ thống chatbot tư vấn bán hàng nội thất theo mô hình multi-tenant kết hợp với kỹ thuật Retrieval-Augmented Generation. Mô hình multi-tenant được sử dụng để nhiều cửa hàng có thể cùng vận hành trên một nền tảng, nhưng mỗi cửa hàng vẫn có dữ liệu sản phẩm, chatbot và cấu hình riêng. Kỹ thuật RAG được sử dụng để chatbot có thể truy xuất thông tin từ knowledge base trước khi tạo câu trả lời, giúp nội dung tư vấn bám sát dữ liệu sản phẩm hơn.

Giải pháp của đề tài được triển khai theo các bước chính sau. Trước hết, hệ thống cho phép cửa hàng quản lý dữ liệu sản phẩm hoặc tài liệu sản phẩm của mình. Các dữ liệu này được xử lý để xây dựng knowledge base riêng cho từng cửa hàng. Khi người dùng đặt câu hỏi, chatbot sẽ xác định ngữ cảnh hội thoại, truy xuất các thông tin liên quan từ knowledge base, sau đó tạo câu trả lời tư vấn dựa trên dữ liệu tìm được.

Đối với bài toán hội thoại nhiều lượt, hệ thống cần lưu lại các thông tin quan trọng trong quá trình trao đổi, chẳng hạn như loại sản phẩm người dùng đang quan tâm, ngân sách, kích thước mong muốn, phong cách hoặc các sản phẩm đã được gợi ý trước đó. Nhờ đó, khi người dùng thay đổi nhu cầu, yêu cầu gợi ý lại hoặc muốn so sánh thêm sản phẩm, chatbot có thể điều chỉnh câu trả lời theo ngữ cảnh mới thay vì xử lý từng câu hỏi một cách rời rạc.

Đối với chức năng so sánh và tham chiếu giá, hệ thống sử dụng dữ liệu sản phẩm tổng hợp từ các cửa hàng tham gia. Dữ liệu này được dùng để hỗ trợ người dùng so sánh nhiều sản phẩm theo các tiêu chí như giá, chất liệu, kích thước, loại sản phẩm, phong cách và cửa hàng cung cấp. Ngoài ra, hệ thống có thể đưa ra thông tin tham khảo về khoảng giá hoặc mặt bằng giá của một nhóm sản phẩm dựa trên dữ liệu hiện có, giúp người dùng có thêm cơ sở trước khi ra quyết định.

Về kiến trúc tổng quát, hệ thống gồm các thành phần chính: giao diện chat cho người dùng cuối, giao diện quản trị cho cửa hàng, backend nghiệp vụ, AI service, cơ sở dữ liệu và knowledge base. Backend nghiệp vụ chịu trách nhiệm quản lý cửa hàng, người dùng, chatbot, dữ liệu sản phẩm, hội thoại và yêu cầu mua hàng. AI service chịu trách nhiệm xử lý câu hỏi, truy xuất tri thức, quản lý ngữ cảnh hội thoại và tạo câu trả lời tư vấn.

Như vậy, định hướng giải pháp của đề tài tập trung vào việc kết nối dữ liệu sản phẩm với quá trình tư vấn bán hàng. Dữ liệu riêng của từng cửa hàng được sử dụng để phục vụ tư vấn theo cửa hàng, trong khi dữ liệu tổng hợp được sử dụng cho nhu cầu so sánh và tham chiếu giá. Cách tiếp cận này giúp hệ thống vừa đáp ứng nhu cầu riêng của từng cửa hàng, vừa hỗ trợ người mua hàng có thêm thông tin khi lựa chọn sản phẩm nội thất.

% Gợi ý bổ sung hình định hướng giải pháp:
% \begin{figure}[H]
%     \centering
%     \includegraphics[width=0.95\textwidth]{figures/solution_orientation.png}
%     \caption{Định hướng giải pháp của hệ thống}
%     \label{fig:solution-orientation}
% \end{figure}

\section{Bố cục đồ án}
\label{section:1.4}
Phần còn lại của báo cáo đồ án tốt nghiệp này được tổ chức như sau.

\textbf{Chương 2} trình bày quá trình khảo sát bài toán và phân tích nhu cầu của các nhóm người dùng trong hệ thống chatbot tư vấn bán hàng nội thất. Nội dung chương tập trung làm rõ bối cảnh sử dụng hệ thống, các đối tượng liên quan như người dùng cuối, cửa hàng và quản trị hệ thống, đồng thời phân tích các yêu cầu chức năng và phi chức năng cần thiết. Chương này cũng trình bày phạm vi dữ liệu được sử dụng trong hệ thống, các tình huống hội thoại tiêu biểu và những tiêu chí cần xem xét khi đánh giá chất lượng tư vấn, so sánh sản phẩm và tham chiếu giá.

\textbf{Chương 3} giới thiệu các công nghệ và cơ sở kỹ thuật được sử dụng để xây dựng hệ thống. Chương này trình bày tổng quan về các thành phần công nghệ chính phục vụ phát triển backend, AI service, cơ sở dữ liệu, giao diện người dùng và môi trường triển khai. Bên cạnh đó, chương cũng giới thiệu các kỹ thuật liên quan đến chatbot, knowledge base, truy xuất thông tin và Retrieval-Augmented Generation ở mức tổng quan, nhằm làm cơ sở cho việc thiết kế và triển khai hệ thống trong các chương tiếp theo.

\textbf{Chương 4} trình bày kết quả thực nghiệm và đánh giá hệ thống. Nội dung chương tập trung vào cách xây dựng dữ liệu mẫu, tập câu hỏi benchmark và các kịch bản kiểm thử phục vụ đánh giá. Trên cơ sở đó, chương phân tích kết quả của hệ thống theo các nhóm tiêu chí chính, gồm chất lượng truy xuất từ knowledge base, chất lượng câu trả lời tư vấn, khả năng xử lý hội thoại nhiều lượt, khả năng gợi ý và so sánh sản phẩm, khả năng tham chiếu giá và khả năng tạo yêu cầu mua hàng. Chương này cũng trình bày các nhận xét về điểm đạt được và những hạn chế còn tồn tại sau quá trình thử nghiệm.

\textbf{Chương 5} trình bày giải pháp xây dựng hệ thống và các đóng góp chính của đồ án. Chương này mô tả cách hệ thống được tổ chức để hỗ trợ nhiều cửa hàng trên cùng một nền tảng, trong đó mỗi cửa hàng có dữ liệu, chatbot và cấu hình riêng. Nội dung chương cũng trình bày các thành phần chính của hệ thống, gồm backend nghiệp vụ, AI service, giao diện quản trị, giao diện chat và các nhóm API phục vụ hội thoại, quản lý knowledge base, quản lý yêu cầu mua hàng, so sánh sản phẩm và tham chiếu giá. Đây là chương tập trung làm rõ sản phẩm phần mềm được xây dựng trong đồ án.

\textbf{Chương 6} tổng kết các nội dung đã thực hiện trong đồ án và đánh giá mức độ hoàn thành so với mục tiêu ban đầu. Chương này trình bày những kết quả chính đã đạt được, các hạn chế còn tồn tại trong phạm vi triển khai và đánh giá, đồng thời đề xuất một số hướng phát triển tiếp theo cho hệ thống. Các hướng phát triển có thể bao gồm mở rộng dữ liệu, cải thiện chất lượng truy xuất và câu trả lời, bổ sung kênh tích hợp, nâng cao khả năng quản trị knowledge base và hoàn thiện hệ thống để tiến gần hơn tới môi trường vận hành thực tế.

\end{document}